\subsection{معماری و پیاده‌سازی}
پیاده‌سازی این ماژول در فایل \lr{Codes/framework/compression\_engine/slim\_compressor.py} با کلاس اصلی \lr{SLiMCompressor} انجام یافته است. مشخصه‌های مهندسی این پیاده‌سازی عبارتند از:
\begin{itemize}
    \item \textbf{سازوکار بازگشت امنیتی (\lr{Fallback Mechanism}):} برای هر لایه، خطای نسبی بازسازی با معیار نرم فروبنیوس $\|\mathbf{W}_0 - \mathbf{W}_{slim}\|_F / \|\mathbf{W}_0\|_F$ اندازه‌گیری می‌شود. چنانچه این خطا از آستانه \lr{max\_error} (پیش‌فرض 0.5) فراتر رود، لایه به وزن اصلی بازگردانده می‌شود.
    \item \textbf{پشتیبانی از لایه‌های خطی و کانوولوشن یک‌بعدی (\lr{Conv1D}):} تشخیص مطمئن لایه‌های خطی بر مبنای نوع لایه و مشخصه \lr{nf} (مورد استفاده در مدل‌های خانواده GPT) بدون تکیه بر حدسیات ابعادی انجام می‌گیرد.
    \item \textbf{محافظت از لایه‌های بحرانی:} لایه‌های حساس مانند طبقه‌بند (\lr{classifier})، سر مدل زبانی (\lr{lm\_head}) و لایه‌های نرمال‌سازی با فراخوانی تابع سراسری \lr{should\_compress\_linear} محافظت می‌شوند.
\end{itemize}

\subsection{محدودیت‌ها و تحلیل شکست}
نقاط ضعف و ریسک‌های معماری این فشرده‌ساز عبارتند از:
\begin{itemize}
    \item \textbf{تداخل خطای کوانتیزاسیون و هرس در لایه‌های عمیق:} در صورتی که طیف مقادیر منفرد ماتریس خطا به کندی افت کند، مرتبه کوچک $r$ در تقریب SVD قادر به جبران اثر کوانتیزاسیون نبوده و خطای نسبی از حد مجاز فراتر می‌رود که منجر به فعال‌سازی حالت بازگشت (\lr{fallback}) و عدم فشرده‌سازی واقعی لایه خواهد شد.
    \item \textbf{سربار حافظه‌ای نگهداری مقادیر SVD:} نگهداری همزمان ماتریس‌های تنک و فاکتورهای کم‌مرتبه $\mathbf{U}_r$ و $\mathbf{V}_r$ در صورتی که مرتبه انتخابی بزرگ باشد، مزیت ذخیره‌سازی را خنثی می‌سازد.
\end{itemize}
